\paragraph{\trMappingCone。}\label{sub:mapping-cones}
设 $c\in\hom_{Tw\scr{A}}(Y_0,Y_1)$ 是一条零次\trCocycle。$c$ 的\trMappingCone{}%
$C=\mathrm{Cone}(c)$ 是如下的\trTwistedComplex：
\begin{equation}\label{eq:mapping-cone-twisted-complex}
\Bigl(C=SY_0\oplus Y_1,\ \delta_C=
\begin{pmatrix}S(\delta_{Y_0}) & 0\\ -S(c) & \delta_{Y_1}\end{pmatrix}\Bigr).
\end{equation}
这里，我们利用\trShiftFunctor{}把 $c$ 变成元素 $S(c)\in\hom^0_{Tw\scr{A}}(SY_0,SY_1)$，
随后以\trTautological{}的方式把该空间等同于 $\hom^1_{Tw\scr{A}}(SY_0,Y_1)$；比较
\eqref{eq:shift-inverse-isomorphism}。若每个 $Y_k$ 都有一个由子\trComplex{}组成的%
\trFiltration{} $F^iY_k$，使得诸\trDifferential{}对此是严格下三角的，那么\trMappingCone{}%
上合成的\trFiltration（先 $F^iSY_0\oplus Y_1$，继而 $0\oplus F^iY_1$）对 $C$ 而言亦然。
广义\trMaurerCartan{}方程恰与如下条件相\trEquivalence：诸 $Y_k$ 各自的方程，连同
$c$ 是 $Tw\scr{A}$ 中的\trCocycle{}这一事实。$C$ 在\trYonedaEmbedding{}下的像可以等同于抽象\trMappingCone{}%
$\scr{C}=\scr{C}one(c)$，只需
\begin{equation*}
\begin{split}
\hom_{Tw\scr{A}}(X,C)&=\hom_{Tw\scr{A}}(X,SY_0)\oplus\hom_{Tw\scr{A}}(X,Y_1)\\
&\cong\hom_{Tw\scr{A}}(X,Y_0)[1]\oplus\hom_{Tw\scr{A}}(X,Y_1)=\scr{C}(X).
\end{split}
\end{equation*}
总而言之，我们已经看到：$Tw\scr{A}$ 含有代表下列运算的特选\trObject{}——\trDirectSum；
与有限维\trGraded\trVectorSpace{}的\trTensorProduct{}（因而特别是\trSuspension；以及在%
\trStrictlyUnital{}的情形，与有限维\trChainComplex{}的\trTensorProduct{}）；而最重要的，
则是\trMappingCone。今后凡与\trTwistedComplex{}打交道，我们自然会使用这些\trObject{}作为代表。

\vkthmsync{26}
\begin{remark}\label{thm:shifted-multiplicity-spaces}
作为另一种方案，也可以只允许一维的\trMultiplicity{}空间来定义\trAdditive{}扩大；若要
更加具体，还可以只允许 $\KK$ 的\trSuspension{}版本。于是，一个\trObject{} $X$ 将写成%
\trDirectSum{} $\bigoplus_i S^{-\sigma^i}X^i$ 的形式，其中 $X^i\in\Ob\scr{A}$ 且
$\sigma^i\in\ZZ$。认定
\begin{equation}\label{eq:shifted-multiplicity-identification}
\hom_{\Sigma\scr{A}}\Bigl(\bigoplus_i S^{-\sigma^i_0}X^i_0,
\bigoplus_j S^{-\sigma^j_1}X^j_1\Bigr)
=\bigoplus_{i,j}\hom_{\scr{A}}(X^i_0,X^j_1)[\sigma^j_1-\sigma^i_0]
\end{equation}
即取显然的\trIsomorphism，并为每个\trCoefficients{} $a^{ji}$ 配备一个附加的正负号
$(-1)^{(|a^{ji}|-1)\sigma^i_0}$（这里的 $|a^{ji}|$ 是 $\hom_{\scr{A}}(X^i_0,X^j_1)$ 中、
在任何\trSuspension{}施加之前的\trDegree）。相对于这一等同，\trComposition{}变为
\begin{equation*}
\mu^d_{\Sigma\scr{A}}(a_d,\ldots,a_1)^{i_d,i_0}=(-1)^{\sigma^{i_0}_0}
\sum_{i_1,\ldots,i_{d-1}}\mu^d_{\scr{A}}\bigl(a_d^{i_d,i_{d-1}},\ldots,a_1^{i_1,i_0}\bigr),
\end{equation*}
其中 $a_k\in\hom_{\Sigma\scr{A}}(X_{k-1},X_k)$，而 $\sigma^{i_0}_0$ 是 $X_0$ 的第 $i_0$
个\trDirectSummand{}上的\trSuspension。\trShiftFunctor{} $S$ 把\trMorphism{}%
$a\in\hom^k_{\Sigma\scr{A}}(X_0,X_1)$ 乘以 $(-1)^{k-1}$。如果接着沿同样的路线去定义一个
$Tw\scr{A}$ 的版本，那么 $C=\mathrm{Cone}(c)$ 上的\trDifferential{}将会是
\begin{equation*}
\delta_C=\begin{pmatrix}\delta_{Y_0} & 0\\ c & \delta_{Y_1}\end{pmatrix};
\end{equation*}
\eqref{eq:mapping-cone-twisted-complex} 中的\trSuspension{}运算与正负号已被吸收进等同
\eqref{eq:shifted-multiplicity-identification} 之中。在\trQuasiequivalence{}的意义下，
允许一般的\trMultiplicity{}空间还是只允许 $\KK$ 的\trSuspension{}版本，当然无关紧要。
更一般的做法有一个优点：容许\trTensorProduct{}运算 $Z\otimes-$ 而无需任何\trBasis{}的
选取（它带来的正负号也更成系统，尽管略为复杂）；尽管如此，二者的总体差别微乎其微。
\end{remark}
